\section{System and Network Security Applications}

\subsection{Multi-Factor Authentication}

Multi-Factor Authentication (MFA) is a familiar (and evocative)
term for Purdue students. 

We require authentication to log into a computer resource (network, device,
application, etc.).
User provides a piece of evidence (factor) for their identity: e.g., a password.
Multi-factor (two-factor) authentication requires two or more pieces of such evidence,
based on the premise that an attacker is unlikely to obtain both. 

Factors can include physical objects (security token via a USB stick, card),
something the user knows (password, PIN), something the user is (biometrics such
as finger print or typing patterns), or location (within a network).

An algorithm is used to create a new
sequence of characters valid for a
single login session.
The generated sequence of characters
is essentially a one-time pad (OTP).

A secret key is shared between the client and server.
Each passcode should be used only once,
some variable to generate a different
OTP each time user wants to log in.
\begin{itemize}
    \item HOTP: HMAC-based one-time pass: variable is a counter.
    \item TOTP: Time-based one-time pass: variable is the current time.
\end{itemize}
TOTP is superior in almost every respect and is widely 
used in industry. 

\subsection{Universal Second Factor}

Shared secret may be transmitted in plaintext between the client and server.
Shared secret may be stored in plaintext on server.
Passcode vulnerable to MITM attack.
Passcode has to be entered manually, prone to error.
Battery operated.
Universal 2nd Factor (U2F) is an open standard that addresses these issues.

U2F device and client (e.g.,
web browser) are controlled
by the user, and relying party
is controlled by the server
(doing authentication).
U2F is based on a
challenge-response protocol.
U2F device holds a public key
pair, and signs a challenge
sent by the server during
authentication.
